\appendices
\section{Implementation and Evaluation Settings}
\label{app:settings}
\subsection{Interpretation Backends}
The studies were executed on September 18, 2026. Jev uses OpenRouter's
Decisions endpoint with identifier typesafe/jev-1.13; responses identify
jev-1.13-20260917 and provider TypeSafe. Four native Choice questions return
the shared intent fields in one request. DeepSeek uses deepseek-v4.1-flash
through OpenRouter, pinned to Together with provider fallback disabled.
It returns a strict four-field JSON object with temperature zero, a
128-token limit, and reasoning disabled; recorded reasoning-token counts
are zero.

Qwen2.5-7B-Instruct, revision a09a35458c70, runs in brain floating point 16
(BF16) on one NVIDIA L40, using Transformers with scaled dot-product
attention (SDPA) and PyTorch 2.2.0+cu121. It uses the same generative prompt
and field contract as DeepSeek, greedy decoding, and a 128-token limit.
The server batches up to four already-pending requests without an intentional
collection delay. The fixed rule parser uses keywords and regular expressions
without tuning on evaluation wording.

\subsection{Admission and Execution Controls}
Both studies use a 2\,s request deadline, four concurrent interpretation
slots, and an admission queue of capacity 32. API calls have a 15\,s timeout,
with no automatic retries or repair calls. Connections are warmed before
measurement. Each condition starts with an empty cache keyed by normalized
full request text and extraction policy. The controller checks it before
queueing and after obtaining a slot; concurrent misses are not coalesced.

The real service has one local macOS worker and one remote Linux worker,
connected through a Secure Shell (SSH) tunnel. The remote allocation provides
six processor cores and the accelerator used by Qwen; OCR uses processors
only. Each node has one nonpreemptive worker and a queue of capacity 16.
Urgent and ordinary jobs have priority values zero and one, respectively;
equal-priority jobs follow enqueue order. The controller verifies the returned
node, tier, priority, and image identity. A timed-out OCR process is terminated
before capacity is released. Transport or protocol errors stop further calls.
The recorded environment lacks a complete processor inventory for exact
cross-site reproduction.

\subsection{Study A Parameters}
The 216 requests use six wording families and two instances of each of the
108 field combinations, with extraction instructions held fixed. Seeds 42,
43, and 44 vary request order and arrival traces across three consecutive
blocks. Each model has 648 semantic measurements on the same 216 texts.
There is one semantic-stage warmup per model per block and no fixed interval
between blocks.

Each service condition has 48 arrivals. Steady traffic has rate 2 requests/s;
bursty traffic alternates 0.5 and 8 requests/s in 20\,s segments. The execution
model has four edge nodes and one cloud node, with payloads of 0.25--2\,MB.
Local, other-edge, and cloud links have propagation delays of 2, 20, and
60\,ms and bandwidths of 1000, 100, and 50\,Mbit/s, respectively; transfer
also includes serialization. Base durations are 40\,ms for counting, 80\,ms
for detection, and 60\,ms for OCR. The high tier multiplies these by 1.8;
node factors are 1.0, 1.1, 1.2, and 1.3 for the edges and 0.65 for the cloud.
Each node has one priority worker. These are exploratory simulation inputs.

\subsection{Study B Parameters and Scoring}
Model conditions run sequentially in randomized order with seed 42 and no
fixed cooldown. Both workers use Tesseract 5.5.3, English recognition,
one OpenMP thread, and single-word mode (engine mode 1, segmentation mode 8).
Standard and high tiers use tessdata\_fast and tessdata\_best, respectively,
with identical weights across nodes. Seed 42 selects 20 training and 80 test
images from sorted IIIT5K identifiers before recognition outcomes are known.

Calibration runs each training image once per node and tier. Median
request--response durations, including transfer, are 80.85/96.26\,ms for
local standard/high and 95.30/139.35\,ms for remote standard/high.
The workload has 48 distinct intents. Repeated-text conditions use eight
descriptions, six supported and two unsupported, each repeated six times.
The test images recur across conditions. Arrival rates match Study A;
eight conditions for four backends produce 32 runs and 1536 arrivals.

OCR comparison uses case-preserving character annotations without a
recognition lexicon, normalizes Unicode to Normalization Form C (NFC), and
strips surrounding whitespace. Quantile summaries use nearest-rank 5th and
95th percentiles and linearly interpolated interquartile ranges. Each
request is evaluated on its recorded timing and output; cached decisions
do not replace service execution.
